\chapter{Facial Bone Fractures}

\section*{Definition and Overview}
Facial bone fractures are breaks in the bones of the face, including the nose, cheekbones, eye sockets, jaws, and the bones around the eyes and forehead. They are usually caused by trauma, such as falls, assaults, sporting injuries, and motor vehicle accidents. Because the face contains the eyes, nose, airways, and important nerves, facial fractures can affect vision, breathing, eating, and appearance. Most facial fractures are treated by specialist teams, and outcomes are generally good with proper care.

\section*{Epidemiology}
Facial fractures are common in trauma and are more frequent in men, particularly in younger adults, with assault, sport, and road crashes the main causes. The nose is the most commonly fractured facial bone, followed by the cheekbone and jaw. In older people, falls are a common cause, and alcohol plays a role in many injuries. Most facial fractures occur alongside other injuries that must be identified first.

\section*{Aetiology and Pathophysiology}
Facial fractures result from direct force to the face.
\begin{itemize}
    \item \textbf{Causes:} falls, assaults, sporting injuries, motor vehicle accidents, and workplace accidents
    \item \textbf{Common fractures:} nasal bone, cheekbone (zygoma), eye socket (orbital), upper and lower jaw, and the bones around the eye
    \item \textbf{Mechanism:} the force breaks the bone and may displace the fragments
    \item \textbf{Associated injuries:} head injury, eye injury, dental injury, and airway problems often accompany facial fractures
    \item \textbf{Effects:} swelling, deformity, pain, and functional problems with breathing, eating, and vision
    \item \textbf{Healing:} fractures heal well when the fragments are aligned, either naturally or with surgery
\end{itemize}

\section*{Clinical Features}
The signs depend on which bones are broken.
\begin{itemize}
    \item Swelling, bruising, and tenderness of the face
    \item Deformity, such as a flattened cheek or a misaligned jaw
    \item A nose that is bent, bleeds, or is blocked
    \item Difficulty opening the mouth or an abnormal bite
    \item Numbness of the cheek, lip, or chin
    \item Double vision, bruising around the eye, or a sunken eye
    \item Bruising behind the ear or around both eyes
    \item Numb or loose teeth
\end{itemize}

\section*{Differential Diagnosis}
Other injuries and conditions can mimic facial fractures.
\begin{itemize}
    \item Soft tissue injury with swelling but no fracture
    \item Dental injuries, including fractured or avulsed teeth
    \item Eye injuries without fractures
    \item Head injury, including concussion and intracranial bleeding
    \item Sinusitis and other causes of facial pain
\end{itemize}

\section*{When to Seek Medical Care}
Anyone with significant facial trauma, deformity, difficulty breathing, double vision, or an abnormal bite should be assessed urgently in an emergency department.

\textbf{Call 000 (or local emergency services) for:}
\begin{itemize}
    \item Difficulty breathing or signs of airway obstruction
    \item Loss of consciousness, confusion, or signs of head injury
    \item Sudden loss of vision or severe double vision
    \item Severe bleeding from the nose or mouth
    \item A facial fracture with spinal injury suspected
\end{itemize}

\section*{Diagnosis and Investigations}
Facial fractures are diagnosed with imaging.
\begin{itemize}
    \item \textbf{Examination:} assessment of the face, eyes, mouth, and bite
    \item \textbf{CT scan:} the key imaging test, showing the fractures in detail
    \item \textbf{Plain X-rays:} used for some nasal and jaw fractures
    \item \textbf{Eye assessment:} for fractures involving the eye socket
    \item \textbf{Dental assessment:} for jaw fractures
    \item \textbf{Assessment of other injuries:} head, neck, and chest
\end{itemize}

\section*{Management}
Management depends on the type and severity of the fracture.
\begin{itemize}
    \item \textbf{Observation:} minor, undisplaced fractures may heal with observation, pain relief, and avoiding pressure
    \item \textbf{Nasal fractures:} sometimes reduced (straightened) under anaesthetic
    \item \textbf{Open reduction and fixation:} surgery to realign the bones and hold them with plates and screws
    \item \textbf{Jaw fractures:} often treated with plates, or with wiring the jaws together
    \item \textbf{Orbital fractures:} surgery if the eye is affected
    \item \textbf{Antibiotics:} for some fractures involving the sinuses or teeth
    \item \textbf{Follow-up:} review of healing, vision, and the bite
\end{itemize}

\section*{Complications}
\begin{itemize}
    \item Infection, particularly of fractures involving the sinuses or teeth
    \item Nerve damage causing permanent numbness or weakness
    \item Double vision or damage to the eye
    \item Malunion, with a change in the shape of the face or bite
    \item Nasal blockage and a crooked nose after healing
    \item Scarring and the need for further surgery
\end{itemize}

\section*{Prognosis and Outcomes}
Most facial fractures heal well, and with proper treatment the face returns to normal or near-normal appearance and function. Early treatment improves outcomes, and serious complications are uncommon with specialist care. Some people are left with lasting numbness, minor asymmetry, or jaw problems, but these are usually manageable.

\section*{Prevention and Self-Care}
\begin{itemize}
    \item Wear seatbelts and helmets
    \item Wear mouthguards and face protection in contact sports
    \item Avoid alcohol and risk-taking behaviour
    \item Use protective equipment at work
    \item Make the home safe to prevent falls, especially for older people
    \item Follow the treatment plan after injury
    \item Attend follow-up and report any changes
\end{itemize}

\section*{Sources and Further Reading}
The following reputable sources provide further information for healthcare students and patients seeking reliable, current advice about facial fractures.
\begin{itemize}
    \item Australian and New Zealand Association of Oral and Maxillofacial Surgeons. \textit{Facial fracture information.} https://www.anzoms.org/
    \item Healthdirect Australia. \textit{Facial injuries.} https://www.healthdirect.gov.au/
    \item NSW Health. \textit{Facial trauma information.} https://www.health.nsw.gov.au/
    \item NHS. \textit{Nasal fractures and facial injuries.} https://www.nhs.uk/conditions/
    \item Mayo Clinic. \textit{Facial fractures.} https://www.mayoclinic.org/diseases-conditions/
    \item MedlinePlus. \textit{Nasal fracture and facial trauma.} https://medlineplus.gov/
\end{itemize}
